\documentclass[12pt,a4paper]{ctexart}

\usepackage[a4paper,margin=2.6cm]{geometry}
\usepackage{amsmath,amssymb,amsthm,mathtools,bm}
\usepackage{physics}
\usepackage{booktabs}
\usepackage{enumitem}

\title{数学分析模拟试卷答案}
\author{}
\date{}

\begin{document}

\maketitle

\begin{center}
\textbf{总分：100 分 \quad 考试时间：120 分钟}
\end{center}

\bigskip

% ==================== 一、填空题 ====================
\section*{一、填空题（每题 5 分，共 20 分）}

\begin{enumerate}[leftmargin=*,label=\arabic*.]
  \item 设 $f(x) = \ln\frac{1+x}{1-x}$，则 $f(x)$ 是\underline{\quad 奇 \quad}函数，
    其定义域为\underline{\quad $(-1,1)$ \quad}。

    \textbf{【解析】}$f(-x) = \ln\frac{1-x}{1+x} = \ln\!\left(\frac{1+x}{1-x}\right)^{-1}
    = -\ln\frac{1+x}{1-x} = -f(x)$，故为奇函数。
    由 $\frac{1+x}{1-x} > 0$ 解得 $-1 < x < 1$，即定义域为 $(-1,1)$。

  \item $\displaystyle\lim_{x\to 0}\frac{e^x - 1 - x}{x^2} = \underline{\quad \dfrac{1}{2} \quad}$。

    \textbf{【解析】}利用 Taylor 展开 $e^x = 1+x+\frac{x^2}{2}+o(x^2)$，
    则 $e^x-1-x = \frac{x^2}{2}+o(x^2)$，故极限为 $\frac{1}{2}$。
    也可用 L'H\^opital 法则：
    $\lim_{x\to 0}\frac{e^x-1-x}{x^2} \overset{H}{=} \lim_{x\to 0}\frac{e^x-1}{2x}
    \overset{H}{=} \lim_{x\to 0}\frac{e^x}{2} = \frac{1}{2}$。

  \item 向量 $\mathbf{a} = (1,2,3)$ 与 $\mathbf{b} = (2,-1,0)$ 的夹角余弦为\underline{\quad $0$ \quad}。

    \textbf{【解析】}$\mathbf{a}\cdot\mathbf{b} = 1\cdot 2 + 2\cdot(-1) + 3\cdot 0 = 2-2+0=0$。
    $\cos\theta = \frac{\mathbf{a}\cdot\mathbf{b}}{|\mathbf{a}||\mathbf{b}|} = 0$，
    故两向量垂直，夹角为 $90^\circ$。

  \item 幂级数 $\displaystyle\sum_{n=1}^{\infty}\frac{2^n}{n}x^n$ 的收敛半径为\underline{\quad $\dfrac{1}{2}$ \quad}。

    \textbf{【解析】}$a_n = \frac{2^n}{n}$，
    $R = \lim_{n\to\infty}\left|\frac{a_n}{a_{n+1}}\right|
    = \lim_{n\to\infty}\frac{2^n/n}{2^{n+1}/(n+1)}
    = \lim_{n\to\infty}\frac{n+1}{2n}
    = \frac{1}{2}$。
\end{enumerate}

% ==================== 二、计算题 ====================
\section*{二、计算题（每题 10 分，共 40 分）}

\begin{enumerate}[leftmargin=*,label=\arabic*.,start=5]
  \item \textbf{（10 分）}设 $y = \left(\frac{\sin x}{x}\right)^{\arctan x}\;(x>0)$，求 $\frac{dy}{dx}$。

    \textbf{解：}取对数 $\ln y = \arctan x \cdot \ln\frac{\sin x}{x}$。
    两边对 $x$ 求导（利用隐函数求导法）：
    \begin{align*}
      \frac{y'}{y} = \frac{1}{1+x^2}\ln\frac{\sin x}{x}
      + \arctan x \cdot \frac{1}{\frac{\sin x}{x}} \cdot
      \frac{x\cos x - \sin x}{x^2}.
    \end{align*}
    又 $\frac{1}{\frac{\sin x}{x}} \cdot \frac{x\cos x - \sin x}{x^2}
    = \frac{x}{\sin x} \cdot \frac{x\cos x - \sin x}{x^2}
    = \frac{x\cos x - \sin x}{x\sin x} = \cot x - \frac{1}{x}$。

    故
    \begin{align*}
      \frac{y'}{y} &= \frac{\ln(\sin x/x)}{1+x^2}
      + \arctan x\left(\cot x - \frac{1}{x}\right). \\[4pt]
      \frac{dy}{dx}
      &= \left(\frac{\sin x}{x}\right)^{\arctan x}
        \left[\frac{\ln(\sin x/x)}{1+x^2}
        + \arctan x\left(\cot x - \frac{1}{x}\right)\right].
    \end{align*}

  \item \textbf{（10 分）}计算定积分 $I = \displaystyle\int_0^{\pi} \frac{x\sin x}{1+\cos^2 x} dx$。

    \textbf{解：}令 $x = \pi - t$，则 $dx = -dt$，$x=0$ 时 $t=\pi$，$x=\pi$ 时 $t=0$。
    $\sin x = \sin(\pi-t) = \sin t$，$\cos x = \cos(\pi-t) = -\cos t$，
    $\cos^2 x = \cos^2 t$。
    \begin{align*}
      I &= \int_0^{\pi} \frac{x\sin x}{1+\cos^2 x} dx
         = \int_{\pi}^0 \frac{(\pi-t)\sin t}{1+\cos^2 t} (-dt) \\
        &= \int_0^{\pi} \frac{(\pi-t)\sin t}{1+\cos^2 t} dt
         = \pi\int_0^{\pi} \frac{\sin t}{1+\cos^2 t} dt - I.
    \end{align*}
    故 $2I = \pi\int_0^{\pi} \frac{\sin t}{1+\cos^2 t} dt$。

    计算 $\int_0^{\pi} \frac{\sin t}{1+\cos^2 t} dt$：
    令 $u = \cos t$，$du = -\sin t\,dt$，$t=0$ 时 $u=1$，$t=\pi$ 时 $u=-1$。
    \[
    \int_0^{\pi} \frac{\sin t}{1+\cos^2 t} dt
    = \int_1^{-1} \frac{-du}{1+u^2}
    = \int_{-1}^1 \frac{du}{1+u^2}
    = \bigl[\arctan u\bigr]_{-1}^1
    = \frac{\pi}{4} - \left(-\frac{\pi}{4}\right)
    = \frac{\pi}{2}.
    \]
    因此 $2I = \pi \cdot \frac{\pi}{2} = \frac{\pi^2}{2}$，即 $I = \frac{\pi^2}{4}$。

  \item \textbf{（10 分）}求微分方程 $y'' - 3y' + 2y = e^x$ 的通解。

    \textbf{解：}先解对应的齐次方程 $y'' - 3y' + 2y = 0$。
    特征方程：$\lambda^2 - 3\lambda + 2 = 0$，
    $(\lambda-1)(\lambda-2) = 0$，$\lambda_1=1$，$\lambda_2=2$。
    齐次通解：$y_h = C_1 e^x + C_2 e^{2x}$。

    非齐次项 $f(x) = e^x$。由于 $\lambda=1$ 是单特征根，
    设特解形式为 $y_p = A x e^x$（乘以 $x$ 一次）。
    则 $y_p' = A e^x + A x e^x = A e^x(1+x)$，
    $y_p'' = A e^x(1+x) + A e^x = A e^x(2+x)$。

    代入原方程：
    \begin{align*}
      A e^x(2+x) - 3A e^x(1+x) + 2A x e^x &= e^x, \\
      A e^x[2+x - 3 - 3x + 2x] &= e^x, \\
      A e^x[-1 + 0] &= e^x, \quad \text{即 } -A e^x = e^x.
    \end{align*}
    故 $A = -1$，$y_p = -x e^x$。

    通解为：$y = y_h + y_p = C_1 e^x + C_2 e^{2x} - x e^x$。

  \item \textbf{（10 分）}计算曲面积分 $\displaystyle\iint_S (x^2+y^2) dS$，其中 $S$ 为圆锥面
    $z = \sqrt{x^2+y^2}$ 介于 $z=0$ 与 $z=1$ 之间的部分。

    \textbf{解：}$S: z = \sqrt{x^2+y^2}$，$D: x^2+y^2 \le 1$（锥面在 $z=1$ 的投影）。
    面积元素：
    \[
    dS = \sqrt{1+z_x^2+z_y^2}\,dxdy,
    \]
    其中 $z_x = \frac{x}{\sqrt{x^2+y^2}} = \frac{x}{z}$，
    $z_y = \frac{y}{\sqrt{x^2+y^2}} = \frac{y}{z}$。
    \[
    \sqrt{1+z_x^2+z_y^2} = \sqrt{1+\frac{x^2+y^2}{x^2+y^2}} = \sqrt{2}.
    \]
    故
    \begin{align*}
      \iint_S (x^2+y^2) dS
      &= \iint_D (x^2+y^2) \cdot \sqrt{2} \,dxdy.
    \end{align*}
    采用极坐标 $x=r\cos\theta$，$y=r\sin\theta$，$dxdy = r dr d\theta$：
    \begin{align*}
      \iint_S (x^2+y^2) dS
      &= \sqrt{2}\int_0^{2\pi} d\theta \int_0^1 r^2 \cdot r dr \\
      &= \sqrt{2}\cdot 2\pi \cdot \left[\frac{r^4}{4}\right]_0^1
       = \sqrt{2}\cdot 2\pi \cdot \frac{1}{4}
       = \frac{\sqrt{2}\pi}{2}.
    \end{align*}
\end{enumerate}

% ==================== 三、证明题 ====================
\section*{三、证明题（每题 15 分，共 30 分）}

\begin{enumerate}[leftmargin=*,label=\arabic*.,start=9]
  \item \textbf{（15 分）}设 $f$ 在 $[0,1]$ 上连续，在 $(0,1)$ 内可导，且 $f(0)=0$，$f(1)=1$。
    证明：
    \begin{enumerate}
      \item 存在 $\xi \in (0,1)$ 使 $f(\xi) = 1-\xi$；
      \item 存在两个不同的点 $\eta_1, \eta_2 \in (0,1)$ 使 $f'(\eta_1)f'(\eta_2)=1$。
    \end{enumerate}

    \textbf{证：(1)} 构造辅助函数 $F(x) = f(x) + x - 1$。
    $F$ 在 $[0,1]$ 上连续。$F(0) = f(0) + 0 - 1 = -1 < 0$，
    $F(1) = f(1) + 1 - 1 = 1 > 0$。
    由零点存在定理，存在 $\xi \in (0,1)$ 使 $F(\xi)=0$，即 $f(\xi) = 1-\xi$。

    \textbf{(2)} 由 (1) 知存在 $\xi \in (0,1)$ 使 $f(\xi) = 1-\xi$。
    在 $[0,\xi]$ 上对 $f$ 应用 Lagrange 中值定理：存在 $\eta_1 \in (0,\xi)$ 使
    \[
    f'(\eta_1) = \frac{f(\xi)-f(0)}{\xi-0} = \frac{1-\xi}{\xi}.
    \]
    在 $[\xi,1]$ 上对 $f$ 应用 Lagrange 中值定理：存在 $\eta_2 \in (\xi,1)$ 使
    \[
    f'(\eta_2) = \frac{f(1)-f(\xi)}{1-\xi} = \frac{1-(1-\xi)}{1-\xi} = \frac{\xi}{1-\xi}.
    \]
    由于 $\eta_1 \in (0,\xi)$，$\eta_2 \in (\xi,1)$，且 $\xi > 0$ 且 $\xi < 1$，
    故 $\eta_1 \neq \eta_2$。两式相乘得：
    \[
    f'(\eta_1) f'(\eta_2) = \frac{1-\xi}{\xi} \cdot \frac{\xi}{1-\xi} = 1.
    \]

  \item \textbf{（15 分）}设 $f(x,y)$ 在全平面上二阶连续可微，且满足 Laplace 方程
    $\frac{\partial^2 f}{\partial x^2} + \frac{\partial^2 f}{\partial y^2} = 0$。
    若 $u(r,\theta) = f(r\cos\theta, r\sin\theta)$，证明：
    \[
    \frac{\partial^2 u}{\partial r^2} + \frac{1}{r}\frac{\partial u}{\partial r}
    + \frac{1}{r^2}\frac{\partial^2 u}{\partial\theta^2} = 0.
    \]

    \textbf{证：}令 $x = r\cos\theta$，$y = r\sin\theta$，则 $u(r,\theta) = f(x,y)$。
    利用多元复合函数的链式法则求偏导数。

    先求一阶偏导数：
    \begin{align*}
      \frac{\partial u}{\partial r}
      &= f_x \frac{\partial x}{\partial r} + f_y \frac{\partial y}{\partial r}
       = f_x \cos\theta + f_y \sin\theta, \\[4pt]
      \frac{\partial u}{\partial\theta}
      &= f_x \frac{\partial x}{\partial\theta} + f_y \frac{\partial y}{\partial\theta}
       = -f_x r\sin\theta + f_y r\cos\theta.
    \end{align*}

    求二阶偏导数。先求 $\frac{\partial^2 u}{\partial r^2}$：
    \begin{align*}
      \frac{\partial^2 u}{\partial r^2}
      &= \frac{\partial}{\partial r}(f_x \cos\theta + f_y \sin\theta) \\
      &= (f_{xx} \cos\theta + f_{xy} \sin\theta)\cos\theta
       + (f_{yx} \cos\theta + f_{yy} \sin\theta)\sin\theta \\
      &= f_{xx}\cos^2\theta + 2f_{xy}\sin\theta\cos\theta + f_{yy}\sin^2\theta.
    \end{align*}
    （其中利用了 $f_{xy} = f_{yx}$，因 $f$ 二阶连续可微。）

    再求 $\frac{\partial^2 u}{\partial\theta^2}$：
    \begin{align*}
      \frac{\partial u}{\partial\theta} &= r(-f_x \sin\theta + f_y \cos\theta), \\
      \frac{\partial^2 u}{\partial\theta^2}
      &= \frac{\partial}{\partial\theta}[r(-f_x \sin\theta + f_y \cos\theta)] \\
      &= r\Big[(-f_{xx}(-r\sin\theta) - f_{xy}(r\cos\theta))\sin\theta
               - f_x \cos\theta \\
      &\qquad   + (f_{yx}(-r\sin\theta) + f_{yy}(r\cos\theta))\cos\theta
               + f_y(-\sin\theta)\Big] \\
      &= r^2\bigl[f_{xx}\sin^2\theta - 2f_{xy}\sin\theta\cos\theta
                 + f_{yy}\cos^2\theta\bigr]
         - r(f_x\cos\theta + f_y\sin\theta).
    \end{align*}
    注意最后一项恰好是 $-r \frac{\partial u}{\partial r}$。

    现在代入要证的表达式左端：
    \begin{align*}
      &\frac{\partial^2 u}{\partial r^2} + \frac{1}{r}\frac{\partial u}{\partial r}
        + \frac{1}{r^2}\frac{\partial^2 u}{\partial\theta^2} \\
      = &\bigl[f_{xx}\cos^2\theta + 2f_{xy}\sin\theta\cos\theta + f_{yy}\sin^2\theta\bigr] \\
        &+ \frac{1}{r}\frac{\partial u}{\partial r}
          + \frac{1}{r^2}\Big[r^2(f_{xx}\sin^2\theta - 2f_{xy}\sin\theta\cos\theta
            + f_{yy}\cos^2\theta) - r\frac{\partial u}{\partial r}\Big] \\
      = &\ f_{xx}(\cos^2\theta+\sin^2\theta) + f_{yy}(\sin^2\theta+\cos^2\theta)
          + 2f_{xy}(\sin\theta\cos\theta - \sin\theta\cos\theta) \\
        &+ \frac{1}{r}\frac{\partial u}{\partial r} - \frac{1}{r}\frac{\partial u}{\partial r} \\[2pt]
      = &\ f_{xx} + f_{yy}.
    \end{align*}
    由条件 $f_{xx} + f_{yy} = 0$（Laplace 方程），故原式 $= 0$。证毕。
\end{enumerate}

% ==================== 四、综合题 ====================
\section*{四、综合题（每题 10 分，共 10 分）}

\begin{enumerate}[leftmargin=*,label=\arabic*.,start=11]
  \item \textbf{（10 分）}设 $f(x) = \sum_{n=0}^{\infty} a_n x^n$ 的收敛半径为 $R > 0$，且 $a_0 = 0$。
    证明：
    \begin{enumerate}
      \item $\int_0^x \frac{f(t)}{t} dt = \sum_{n=1}^{\infty} \frac{a_n}{n} x^n$，$|x| < R$；
      \item 利用上述结果求 $\displaystyle\sum_{n=1}^{\infty}\frac{1}{n 2^n}$ 的和。
    \end{enumerate}

    \textbf{证：(1)} 由条件 $a_0=0$，有 $f(t) = \sum_{n=1}^{\infty} a_n t^n$，$|t| < R$。
    则 $\frac{f(t)}{t} = \sum_{n=1}^{\infty} a_n t^{n-1}$。
    幂级数在收敛区间内可逐项积分。对 $|x| < R$，取从 $0$ 到 $x$ 的积分路径包含在收敛区间内：
    \begin{align*}
      \int_0^x \frac{f(t)}{t} dt
      &= \int_0^x \sum_{n=1}^{\infty} a_n t^{n-1} dt
       = \sum_{n=1}^{\infty} a_n \int_0^x t^{n-1} dt
       = \sum_{n=1}^{\infty} a_n \left[\frac{t^n}{n}\right]_0^x \\
      &= \sum_{n=1}^{\infty} \frac{a_n}{n} x^n.
    \end{align*}
    逐项积分的合理性由幂级数在收敛区间内闭一致收敛保证。

    \textbf{(2)} 取 $f(x) = \frac{x}{1-x}$，其收敛半径为 $R=1$，
    $f(x) = \sum_{n=1}^{\infty} x^n = x + x^2 + x^3 + \cdots$，
    即 $a_n = 1$（$n \ge 1$），$a_0=0$ 成立。

    由 (1) 的结果：
    \[
    \int_0^x \frac{f(t)}{t} dt
    = \int_0^x \frac{1}{1-t} dt
    = -\ln(1-x)
    = \sum_{n=1}^{\infty} \frac{1}{n} x^n.
    \]
    取 $x = \frac{1}{2}$（满足 $|x| < 1$）：
    \[
    \sum_{n=1}^{\infty}\frac{1}{n 2^n} = -\ln\!\left(1-\frac{1}{2}\right)
    = -\ln\frac{1}{2}
    = \ln 2.
    \]
    故 $\displaystyle\sum_{n=1}^{\infty}\frac{1}{n 2^n} = \ln 2$。
\end{enumerate}

\end{document}
